\chapter*{Conclusion générale}
\addcontentsline{toc}{chapter}{Conclusion générale}
\markboth{Conclusion générale}{Conclusion générale}

Le projet relie une expérience de classification sur RT-IoT2022 à une
application d'investigation. La préparation contrôle le fichier source et
retire les conflits d'étiquettes ainsi que les doublons. La comparaison de
quatre familles de modèles retient un gradient boosting par histogrammes pour
la détection binaire et une forêt aléatoire pour les familles d'attaque.
L'API enregistre leurs décisions, et l'interface permet de retrouver les
observations, les versions des modèles et les qualifications de l'analyste.

La cascade atteint une macro-F1 de 0,9520 sur le test de la graine~42 et une
moyenne de 0,9686 sur les trois graines. Les résultats sont moins stables pour
les classes rares, dont certaines ne comptent que six ou sept observations de
test. Les partitions aléatoires d'un même corpus ne renseignent pas encore
sur le comportement du modèle dans un autre environnement.

Sur le plan applicatif, la file persistante, les reprises et la publication
après transaction permettent de suivre le traitement d'une observation jusqu'à
sa consultation. La répétition locale documente ces parcours et les scénarios
de signature Suricata. La traçabilité reste toutefois incomplète pour une
reproduction exacte de l'entraînement : le rapport signale des modifications
non commitées, et le matériel des mesures de latence n'est pas détaillé.

La prochaine expérience devrait partir d'une collecte NFStream étiquetée,
avec les paramètres d'extraction, les sessions et les horodatages conservés.
Elle permettrait de séparer les sessions d'entraînement et de test, puis de
mesurer les attaques manquées et les fausses alertes sur un trafic représentatif.
L'archivage du code exact et des caractéristiques de l'hôte compléterait les
fichiers déjà conservés. Ces éléments permettraient de juger le transfert du
modèle au-delà de RT-IoT2022.
